\section*{Coda: compression is a scientific choice}
\addcontentsline{toc}{section}{Coda: compression is a scientific choice}

An operator graph supports several non-equivalent compressions.  Support thresholding asks which
edges survive.  Low-rank approximation asks which substitution directions survive on each edge.
Pair truncation asks which higher-order components survive.  Scalarization asks which single
summary of an edge is retained.  A model can be economical in one sense and expensive in another.

This separation changes how graph comparisons should be reported.  Edge counts measure support;
categorical capacity additionally depends on the maps carried by those edges.  Scalar spectra
compare one compression of operator action, and contact maps evaluate one structural endpoint.
Shared structure may therefore remain in the full operators even when one scalar compression
performs poorly.

The same separation controls the physical interpretation.  A higher-order ANOVA component
certifies irreducibility to lower-order functions in the declared variables and reference measure.
The origin of that irreducibility remains an empirical question.  It may reflect a genuine
many-body term in an effective Hamiltonian, a pair interaction modified by an unobserved solvent or
conformational coordinate, the marginalization of latent states, or nonlinear composition inside
a predictor.  Calling it a three-body or many-body physical force requires the stronger claim that
the decomposed function is itself a calibrated energy on a justified coarse-grained state space.

Low rank has a similarly limited physical license.  A dominant singular direction is a collective
channel of categorical response.  It becomes a normal mode only after a self-adjoint physical
operator, an inner product with the appropriate metric, and a dynamics or quadratic energy have
been specified.  The scalar norm of an edge is weaker still: it measures the magnitude of action
after gauge projection.  Sign, reciprocity, and work around a closed cycle live in additional
coordinates.  Thus
support, rank, order, and norm are possible signatures of physical organization, not physical
observables by definition.

So far the geometry has been linear: tables are projected into centered spaces and compared by
weighted norms.  Probabilities add a nonlinear constraint.  A perturbed pair law must remain
positive and preserve its declared marginals at finite displacement as well as to first order.  The next chapter
shows how the linear gauge arises as the tangent geometry of that probability-exact problem.
